\section{代表性阈值的量级参考}
\label{sec:sup:threshold-scale}
正文西安算例采用 $N=13{,}163{,}000$\cite{XianStat2021}、$\eta=0.002N=26{,}326$。本节保留其比例情景的参考依据，并与上一节的医疗需求充分条件区分。

陈胤孜等\cite{Chen2021ICUForecast}基于既往资料预测，2021 年全国每10万人综合 ICU 床位约为 $4.37$ 张；这一数值为全国预测量，而非西安同期可用容量的实测值。Angulo 等\cite[补充材料 S4.1]{Angulo2021}对早期 COVID-19 采用需要重症监护的感染者比例名义值 $0.60\times0.15\times0.28=0.0252$。将前者作为人均容量参考、不扣除既有占用，并将后者移作参考系数 $\rho_{\mathrm{ICU}}$，按比例情景
\[
\eta=\frac{C}{\rho_{\mathrm{ICU}}},\qquad
\theta=\frac{C/N}{\rho_{\mathrm{ICU}}}
\]
得到
\[
\theta=\frac{4.37/10^5}{0.0252}\approx1.73\times10^{-3},
\qquad \eta\approx2.28\times10^4.
\]
正文在这一量级附近取 $\theta=0.002$ 作代表性设计情景，而非对上式作保守取整。按相同参考系数反算，它对应每10万人 $5.04$ 张、全市约 $663$ 张的名义容量。

上述比例移植用于说明阈值量级，不是当地医疗容量的临床校准。以确诊病例或全部感染者为分母的 ICU 入院风险可能不同\cite{Pung2023Undetected}；此处的参考系数未依据共同需求函数和初始负担估计，也不等同于第\ref{sec:sup:medical-link}节的 $pL\beta c_0S_0/N$。两节分别给出一个历史比例情景和一个附加医疗假设下的充分条件，现有算例采用前者附近的设计阈值开展控制比较。

